\documentclass[11pt]{article}
\usepackage{mathblatt}
\usepackage{adjustbox,varwidth}
% gesetzt von werkzeuge/setzer.py aus dem Lernweg (L6-A · L6-B · L6-C · L6-D · L6-T) und den Bankzeilen; Muster T6B.
% Präambel des Setzers (werkzeuge/setzer.py), wortgleich aus dem Hand-Satz T6B (bau/proben/2026-10-09/kathete-berechnen), dort aus K4W.
\makeatletter
\setlength{\mbpfbe}{0mm}% keine Punkte (Bauregel 6.4)
% eingerückt (Bauregel 4.7, 6.6): \gEin Einzug, \gSchrift eine Stufe kleiner
\newlength{\gEin}\setlength{\gEin}{0pt}
\let\gSchrift\relax
\renewcommand{\pfkreuzzeile}[1]{\par\noindent\hangindent3.2em\hangafter1\hspace*{1.6em}$\square$\ #1\par}
\newcommand{\gkreuz}[1]{$\square$\ #1\hspace{2em}}
% Bauregel 6.7: links, was man ansieht, rechts, was man tut; Spaltengrenze überall an derselben Stelle
\newsavebox{\gbox}\newlength{\gLinks}
\newcommand{\gzwei}[2]{\par\smallskip\noindent\sbox{\gbox}{#1}%
  \setlength{\gLinks}{\dimexpr0.5\textwidth-\gEin-\mbpfnr\relax}%
  \ifdim\wd\gbox>\dimexpr\gLinks-4mm\relax
    \usebox{\gbox}\par\smallskip\noindent\begin{minipage}{\linewidth}#2\end{minipage}%
  \else
    \begin{minipage}[t]{\gLinks}\vspace{0pt}\usebox{\gbox}\end{minipage}%
    \begin{minipage}[t]{\dimexpr\linewidth-\gLinks\relax}\vspace{0pt}\raggedright #2\end{minipage}%
  \fi\par}
\RenewDocumentEnvironment{pfaufg}{m m m}{%
  \begin{lrbox}{\mbaufgabebox}\begin{minipage}[t]{\linewidth}%
  \gSchrift\noindent\hspace*{\gEin}\mbpfrandmarke{#1}%
  \begin{minipage}[t]{\mbpfnr}\strut\textbf{#2}\end{minipage}%
  \begin{minipage}[t]{\dimexpr\linewidth-\gEin-\mbpfnr\relax}\raggedright\strut\ignorespaces}%
 {\end{minipage}%
  \end{minipage}\end{lrbox}%
  \setlength{\mbaufgabehoehe}{\dimexpr\ht\mbaufgabebox+\dp\mbaufgabebox\relax}%
  \par\addvspace{7pt}\Needspace*{\mbaufgabehoehe}%
  \noindent\usebox{\mbaufgabebox}\par\addvspace{7pt}}
\makeatother
% Rechenraum als Karo (Bauregel 6.9): n Rechenschritte = n mal 10 mm
\newcommand{\karo}[1]{\par\smallskip\noindent\begin{tikzpicture}
  \clip (0,0) rectangle ({\linewidth-1mm},{#1*10mm});
  \draw[black!16,line width=0.3pt,step=5mm] (0,0) grid ({\linewidth},{#1*10mm});
  \end{tikzpicture}\par}
\newcommand{\kk}[1]{$\square$\,#1\hspace{0.9em}}
\newcommand{\linie}[1]{\,{\color{black!45}\rule[-1pt]{#1}{0.4pt}}\,}
% Zeile am Ende jedes Durchgangs: Originale klein grau, darüber; dann Vorgänger/Nachfolger (Bauregel 3.4, 6.5)
\newcommand{\blattende}[2]{\par\vfill\noindent\begin{minipage}{\linewidth}{\scriptsize\color{mbgrau}Originale: #1\par}\smallskip
{\footnotesize #2\par}\end{minipage}}
\newcommand{\pkt}{\textperiodcentered{}}

\newcommand{\kennung}{BGG}
\newcommand{\grau}[1]{{\color{mbgrau}#1}}

\begin{document}
\pfheftstil
\fancyfoot[C]{\parbox[t]{\textwidth}{\mbfhbox\par\vspace{1.5mm}{\footnotesize\color{mbgrau}{\scriptsize \kennung}\hfill\thepage}}}
\pfheftkopf{Terme aufstellen}{Klasse 7}
% L6-A: Rechenwörter in Terme übersetzen
\begin{pfaufg}{}{1.}{}
\grau{a)\ Das Doppelte einer Zahl $x$ wird um 7 vermindert. Schreibe dazu einen Term.\par\smallskip Was geschieht zuerst mit $x$? Es wird verdoppelt: \quad $2x$\par „um 7 vermindert“ heißt minus 7, hinten angehängt: \quad $2x - 7$\par Probe mit einer Zahl, z.\,B. $x = 5$: Das Doppelte ist $10$, um 7 vermindert $3$. \ Term: $2 \cdot 5 - 7 = 3$. Passt.}
\par\medskip
b)\ Schreibe als Term. \par\vspace{3pt} a)~das Fünffache von $x$ \qquad b)~$x$ um 8 vergrößert \qquad c)~$x$ um 6 vermindert \par\vspace{3pt} d)~die Hälfte von $x$ \qquad e)~10 vermindert um $x$ \qquad f)~das Produkt aus 7 und $x$\karo{3}
\fusshilfe{\mbox{1:~b) $5x$;\ $x + 8$;\ $x - 6$;\ $x : 2$;\ $10 - x$;\ $7x$}}
\end{pfaufg}

% L6-A: Rechenwörter in Terme übersetzen
\begin{pfaufg}{}{2.}{}
Das Siebenfache einer Zahl $x$ wird um 3 vermindert. Kreuze den passenden Term an.\par\smallskip \gkreuz{$7 \cdot (x - 3)$}\gkreuz{$3 - 7x$}\par \gkreuz{$7x - 3$}\gkreuz{$7x - 3x$}
\fusshilfe{\mbox{2:~$7x - 3$}}
\end{pfaufg}

% L6-A: Rechenwörter in Terme übersetzen
\begin{pfaufg}{}{3.}{}
Schreibe als Term. \par\vspace{3pt} a)~Das Dreifache einer Zahl $x$ wird um 4 vergrößert. \par b)~Das Vierfache von $x$ wird um 9 vermindert. \par c)~20 wird um das Doppelte von $x$ vermindert. \par d)~Die Hälfte von $x$ wird um 1 vergrößert.\karo{3}
\fusshilfe{\mbox{3:~$3x + 4$;\ $4x - 9$;\ $20 - 2x$;\ $x : 2 + 1$}}
\end{pfaufg}

% L6-A: Rechenwörter in Terme übersetzen
\begin{pfaufg}{}{4.}{}
Lena ist $x$ Jahre alt. Ihr Bruder Tim ist 4 Jahre älter als Lena. Ihre Mutter ist dreimal so alt wie Lena. \par a)~Schreibe je einen Term für das Alter von Tim und von der Mutter. \par b)~Lena ist 12 Jahre alt. Ist die Mutter älter als 35 Jahre?\karo{3}
\fusshilfe{\mbox{4:~a) $x + 4$; \ $3x$ \quad b) $36$ Jahre, ja}}
\end{pfaufg}

% L6-B: Klammer, wenn ein Ergebnis weiterverarbeitet wird
\begin{pfaufg}{}{5.}{}
\grau{a)\ Die Differenz aus dem Vierfachen einer Zahl $x$ und 3 wird verdoppelt. Schreibe dazu einen Term.\par\smallskip Das Vierfache von $x$: \quad $4x$\par Die Differenz aus $4x$ und 3: \quad $4x - 3$\par Diese Differenz wird verdoppelt – als Ganzes, also in Klammern: \quad $2 \cdot (4x - 3)$\par Probe mit $x = 2$: Vierfaches $8$, minus 3 ergibt $5$, verdoppelt $10$. \ Term: $2 \cdot (4 \cdot 2 - 3) = 2 \cdot 5 = 10$. Passt.}
\par\medskip
b)\ Schreibe als Term. \par a)~Die Summe aus $x$ und 5 wird verdreifacht. \par b)~Die Differenz aus $x$ und 2 wird mit 4 multipliziert. \par c)~Die Summe aus $x$ und 8 wird halbiert.\karo{3}
\fusshilfe{\mbox{5:~b) $3 \cdot (x + 5)$;\ $4 \cdot (x - 2)$;\ $(x + 8) : 2$}}
\end{pfaufg}

% L6-B: Klammer, wenn ein Ergebnis weiterverarbeitet wird
\begin{pfaufg}{}{6.}{}
Schreibe als Term. Nicht jeder braucht eine Klammer. \par a)~Verdopple eine Zahl $x$, addiere 5 und verdreifache das Ergebnis. \par b)~Die Summe aus dem Dreifachen von $x$ und 4 wird halbiert. \par c)~Vervierfache eine Zahl $x$ und addiere dann 6.\karo{3}
\fusshilfe{\mbox{6:~$3 \cdot (2x + 5)$;\ $(3x + 4) : 2$;\ $4x + 6$}}
\end{pfaufg}

% L6-B: Klammer, wenn ein Ergebnis weiterverarbeitet wird
\begin{pfaufg}{}{7.}{}
Die Differenz aus dem Fünffachen einer Zahl $x$ und 4 wird verdreifacht. Kreuze den passenden Term an.\par\smallskip \gkreuz{$(5x - 4) : 3$}\gkreuz{$5x : 4 \cdot 3$}\par \gkreuz{$3 \cdot (5x - 4)$}\gkreuz{$5x - 4 \cdot 3$}
\fusshilfe{\mbox{7:~$3 \cdot (5x - 4)$}}
\end{pfaufg}

% L6-B: Klammer, wenn ein Ergebnis weiterverarbeitet wird
\begin{pfaufg}{}{8.}{}
Drei Freunde gehen ins Kino. Jeder kauft ein Ticket für $x$ Euro und eine Tüte Popcorn für 5 Euro. Ein Ticket kostet 9 Euro. Reichen den dreien 40 Euro? \par Kreuze zuerst den Term an, der den Gesamtpreis angibt, und rechne dann damit.\par\smallskip \gkreuz{$3x + 5$}\gkreuz{$3 \cdot (x + 5)$}\par \gkreuz{$x + 3 \cdot 5$}\gkreuz{$3 + x + 5$}
\fusshilfe{\mbox{8:~$3 \cdot (x + 5)$; \ $42$ Euro, nein}}
\end{pfaufg}

% L6-C: Term zur Figur: Umfang und Flächeninhalt
\begin{pfaufg}{}{9.}{}
\gzwei{\begin{tikzpicture}[scale=0.45,font=\small]\draw[line width=0.8pt] (0,0) rectangle (8,4);\draw[line width=0.8pt] (5,0) -- (5,4);\node[below] at (2.5,0) {$x$};\node[below] at (6.5,0) {$3$};\node[left] at (0,2.0) {$4$};\end{tikzpicture}}{\grau{a)\ Das Rechteck besteht aus zwei Teilen (Maße in cm). Schreibe Terme für den Umfang und den Flächeninhalt des ganzen Rechtecks.\par\smallskip Ganze Länge: \ $x + 3$; \ Breite: \ $4$\par Umfang, einmal rundherum: \ $u = x + 3 + 4 + x + 3 + 4 = 2x + 14$\par Fläche als Ganzes, Länge mal Breite: \ $A = 4 \cdot (x + 3)$\par Fläche als Summe der Teile: \ $A = 4 \cdot x + 4 \cdot 3 = 4x + 12$. \ Beide Terme sind gleich, Probe mit $x = 5$: $4 \cdot 8 = 32$ \ und \ $20 + 12 = 32$.}}
\par\medskip
\gzwei{\adjustbox{max width=\linewidth}{\begin{varwidth}{50cm}\begin{tikzpicture}[scale=0.4,font=\small]\draw[line width=0.8pt] (0,0) rectangle (4,4);\node[below] at (2.0,0) {$a$};\end{tikzpicture}\qquad \begin{tikzpicture}[scale=0.4,font=\small]\draw[line width=0.8pt] (0,0) rectangle (6,3);\node[below] at (3.0,0) {$x$};\node[left] at (0,1.5) {$5$};\end{tikzpicture}\qquad \begin{tikzpicture}[scale=0.4,font=\small]\draw[line width=0.8pt] (0,0) -- (5,0) -- (2.5,4.33) -- cycle;\node[below] at (2.5,0) {$s$};\end{tikzpicture}\end{varwidth}}}{%
b)\ Schreibe einen Term für den Umfang (Maße in cm). \par a)~Quadrat \qquad b)~Rechteck, dazu auch den Flächeninhalt \qquad c)~Dreieck mit drei gleich langen Seiten\karo{3}}
\fusshilfe{\mbox{9:~b) $4a$;\ $2x + 10$ und $5x$;\ $3s$}}
\end{pfaufg}

% L6-C: Term zur Figur: Umfang und Flächeninhalt
\begin{pfaufg}{}{10.}{}
\gzwei{\begin{tikzpicture}[scale=0.45,font=\small]\draw[line width=0.8pt] (0,0) rectangle (6,4.0);\draw[line width=0.8pt] (0,2.5) -- (6,2.5);\node[above] at (3.0,4.0) {$a$};\node[right] at (6,3.25) {$b$};\node[right] at (6,1.25) {$c$};\end{tikzpicture}}{%
Das Rechteck ist in zwei Streifen geteilt. Welcher Term gibt den Flächeninhalt des ganzen Rechtecks an? Kreuze an.\par\smallskip \gkreuz{$a \cdot b + c$}\gkreuz{$a \cdot (b + c)$}\par \gkreuz{$a + b \cdot c$}\gkreuz{$a \cdot b \cdot c$}}
\fusshilfe{\mbox{10:~$a \cdot (b + c)$}}
\end{pfaufg}

% L6-C: Term zur Figur: Umfang und Flächeninhalt
\begin{pfaufg}{}{11.}{}
\gzwei{\begin{tikzpicture}[scale=0.4,font=\small]\draw[line width=0.8pt] (0,0) rectangle (9,3);\draw[line width=0.8pt] (6,0) -- (6,3);\node[below] at (3.0,0) {$6$};\node[below] at (7.5,0) {$3$};\node[left] at (0,1.5) {$y$};\end{tikzpicture}}{%
Das Rechteck besteht aus zwei Teilen (Maße in m). \par a)~Schreibe zwei verschiedene Terme für den Flächeninhalt. \par b)~Schreibe einen Term für den Umfang.\karo{3}}
\fusshilfe{\mbox{11:~a) $y \cdot (6 + 3)$ und $6y + 3y$ \quad b) $2y + 18$}}
\end{pfaufg}

% L6-C: Term zur Figur: Umfang und Flächeninhalt
\begin{pfaufg}{}{12.}{}
Ein rechteckiges Beet ist $x$ Meter lang und 3 Meter breit. Es bekommt rundherum einen Zaun. \par a)~Skizziere das Beet und schreibe einen Term für die Länge des Zauns. \par b)~Das Beet soll 6 Meter lang werden. Reichen 20 Meter Zaun? \par c)~Wie groß ist dann seine Fläche?\karo{3}
\fusshilfe{\mbox{12:~a) $2x + 6$ \quad b) $18$\,m, ja \quad c) $18$\,m²}}
\end{pfaufg}

% L6-D: Terme in Sachaufgaben
\begin{pfaufg}{}{13.}{}
\grau{a)\ Ein E-Scooter kostet 1 Euro zum Entsperren und 0,20 Euro je Minute. Schreibe einen Term für den Preis einer Fahrt. Was kostet eine Fahrt von 15 Minuten?\par\smallskip Variable festlegen: \ $x$ ist die Zahl der Minuten.\par Fest: 1 Euro. \ Je Minute 0,20 Euro, bei $x$ Minuten: \ $0{,}20 \cdot x$\par Term für den Preis in Euro: \quad $1 + 0{,}20 \cdot x$\par Umgekehrt gelesen: \ 1 ist der feste Betrag, 0,20 der Preis je Minute.\par Für 15 Minuten einsetzen: \ $1 + 0{,}20 \cdot 15 = 1 + 3 = 4$. \ Die Fahrt kostet 4 Euro.}
\par\medskip
b)\ Beim Bowling zahlt man 4 Euro Schuhmiete und 5 Euro je Spiel. Schreibe einen Term für den Preis bei $n$ Spielen. Was kosten 3 Spiele?\karo{3}
\fusshilfe{\mbox{13:~b) $4 + 5n$; \ $19$ Euro}}
\end{pfaufg}

% L6-D: Terme in Sachaufgaben
\begin{pfaufg}{}{14.}{}
Ein Eis kostet $a$ Euro, ein Getränk $b$ Euro. Familie Kaya kauft 4 Eis und 2 Getränke. \par a)~Schreibe einen Term für den Gesamtpreis. \par b)~Ein Eis kostet 1,50 Euro, ein Getränk 2 Euro. Wie viel zahlt die Familie?\karo{3}
\fusshilfe{\mbox{14:~a) $4a + 2b$ \quad b) $10$ Euro}}
\end{pfaufg}

% L6-D: Terme in Sachaufgaben
\begin{pfaufg}{}{15.}{}
Ein Parkhaus berechnet den Preis in Euro mit dem Term $2 + 1{,}50 \cdot x$. \par a)~Was bedeuten die Zahl 2, die Zahl 1,50 und die Variable $x$? \par b)~Was kosten 4 Stunden?\karo{3}
\fusshilfe{\mbox{15:~a) fest 2 Euro, je Stunde 1,50 Euro, $x$ Stunden \quad b) $8$ Euro}}
\end{pfaufg}

% L6-D: Terme in Sachaufgaben
\begin{pfaufg}{}{16.}{}
\gzwei{\adjustbox{max width=\linewidth}{\begin{varwidth}{50cm}\wertetabelle{h}{\text{Fluss}}{2,4,6}\par\wertetabelle{h}{\text{See}}{2,4,6}\end{varwidth}}}{%
Kanuverleih Fluss nimmt 10 Euro Grundgebühr und 2 Euro je Stunde. Kanuverleih See nimmt 4 Euro je Stunde ohne Grundgebühr. \par a)~Schreibe für beide einen Term für die Kosten nach $h$ Stunden. \par b)~Trage die Kosten für 2, 4 und 6 Stunden in die Tabelle ein. \par c)~Familie Roth möchte ein Kanu 6 Stunden leihen. Welcher Verleih ist billiger, und um wie viel?\karo{3}}
\fusshilfe{\mbox{16:~Fluss ($22$ Euro statt $24$ Euro), um $2$ Euro billiger}}
\end{pfaufg}

% L6-T: Probetest
\begin{pfaufg}{}{17.}{}
Das Achtfache einer Zahl $x$ wird um 6 vermindert. Kreuze den passenden Term an.\par\smallskip \gkreuz{$8 \cdot (x - 6)$}\gkreuz{$6 - 8x$}\par \gkreuz{$8x - 6$}\gkreuz{$8x - 6x$}
\fusshilfe{\mbox{17:~$8x - 6$}}
\end{pfaufg}

% L6-T: Probetest
\begin{pfaufg}{}{18.}{}
Die Summe aus dem Doppelten einer Zahl $x$ und 7 wird halbiert. Kreuze den passenden Term an.\par\smallskip \gkreuz{$2x + 7 : 2$}\gkreuz{$(2x + 7) : 2$}\par \gkreuz{$2 \cdot (x + 7) : 2$}\gkreuz{$2x : 7 \cdot 2$}
\fusshilfe{\mbox{18:~$(2x + 7) : 2$}}
\end{pfaufg}

% L6-T: Probetest
\begin{pfaufg}{}{19.}{}
\gzwei{\begin{tikzpicture}[scale=0.35,font=\small]\draw[line width=0.8pt] (0,0) rectangle (6,5);\draw[line width=0.8pt] (4,0) -- (4,5);\node[below] at (2.0,0) {$x$};\node[below] at (5.0,0) {$2$};\node[left] at (0,2.5) {$5$};\end{tikzpicture}}{%
Das Rechteck besteht aus zwei Teilen (Maße in cm). Gib einen Term für den Flächeninhalt und einen für den Umfang an.\karo{3}}
\fusshilfe{\mbox{19:~$A = 5 \cdot (x + 2)$; \ $u = 2x + 14$}}
\end{pfaufg}

% L6-T: Probetest
\begin{pfaufg}{}{20.}{}
Schreibe als Term: Subtrahiere von einer Zahl $x$ die Zahl 3 und multipliziere die Differenz mit 6.\karo{3}
\fusshilfe{\mbox{20:~$6 \cdot (x - 3)$}}
\end{pfaufg}

% L6-T: Probetest
\begin{pfaufg}{}{21.}{}
Ein Bus hat $s$ Sitzplätze und $t$ Stehplätze. Zum Stadion fahren 5 solche Busse. \par a)~Schreibe einen Term für die Zahl der Fahrgäste, die höchstens mitfahren können. \par b)~Wie viele sind es für $s = 40$ und $t = 30$?\karo{3}
\fusshilfe{\mbox{21:~a) $5 \cdot (s + t)$ \quad b) $350$}}
\end{pfaufg}

% L6-T: Probetest
\begin{pfaufg}{}{22.}{}
\gzwei{\adjustbox{max width=\linewidth}{\begin{varwidth}{50cm}\wertetabelle{m}{\text{Fit}}{1,2,3,4,5}\par\wertetabelle{m}{\text{Top}}{1,2,3,4,5}\end{varwidth}}}{%
Studio Fit nimmt 20 Euro Aufnahmegebühr und 15 Euro je Monat. Studio Top nimmt 25 Euro je Monat ohne Aufnahmegebühr. \par a)~Schreibe für beide einen Term für die Kosten nach $m$ Monaten. \par b)~Trage die Kosten in die Tabelle ein. \par c)~Ab welchem Monat ist Fit günstiger?\karo{3}}
\fusshilfe{\mbox{22:~Fit $20 + 15m$, Top $25m$, ab dem 3}}
\end{pfaufg}

\par\vfill\noindent{\footnotesize Vorher: Termwerte berechnen\quad\textbullet\quad Weiter: Terme zusammenfassen\par}
\end{document}
